\section{工艺参数分析}\label{sec:process}

\subsection{情景设置}

以问题二、三的模型为基准，分别改变四个工艺参数，其余条件不变：
\begin{enumerate}[label=(\arabic*),leftmargin=2.2em,itemsep=1pt]
\item 恒温段温度 $T_{\mathrm{set}}$：45--$\SI{70}{\celsius}$，升温段的 $\Tair-\SI{28}{\celsius}$ 按比例缩放。$\SI{40}{\celsius}$ 时含湿量 0.05 已超过饱和含湿量（约 0.049），烘房会结露，故不计入。
\item 恒温段含水浓度 $C_{\mathrm{set}}$：0.02--0.07，前 $\SI{4}{h}$ 的增量按升温进度过渡，初值不变。
\item 药材半径 $R_0$：1.0--$\SI{2.5}{cm}$。
\item 风速：由 Chilton--Colburn 类比\cite{bergman}，$h$ 与 $h_m$ 随风速按相同幂次变化，以同一倍率 0.5--2 同时缩放。
\end{enumerate}
情景计算取 $N=400$（与 $N=800$ 的烘干时长相差 $\SI{1e-4}{h}$）。每个情景同时用题设模型和第~\ref{sec:evap}~节的能量一致模型计算，后者计入了蒸发吸热，结果见图~\ref{fig:process} 和表~\ref{tab:map}。

\begin{figure}[htbp]\centering
\includegraphics[width=\textwidth]{fig_process.pdf}
\caption{工艺参数对烘干时长的影响。(a) 恒温段温度，阴影为药典低温干燥上限 $\SI{60}{\celsius}$ 以上；(b) 恒温段含水浓度；(c) 药材半径（双对数坐标）；(d) 风速。实线为题设模型，虚线为能量一致模型。}
\label{fig:process}
\end{figure}

\begin{table}[htbp]\centering\small
\caption{不同半径与恒温段温度下的烘干时长（题设模型，h）}\label{tab:map}
\begin{tabular}{ccccccc}
\toprule
 & \multicolumn{6}{c}{恒温段温度/\si{\celsius}} \\
\cmidrule(lr){2-7}
初始半径/cm & 45 & 50 & 55 & 60 & 65 & 70 \\
\midrule
1.0 & 18.79 & 16.29 & 14.30 & 12.72 & 11.45 & 10.44 \\
1.5 & 39.46 & 33.67 & 29.04 & 25.34 & 22.37 & 19.96 \\
2.0 & 67.94 & 57.47 & 49.11 & 42.39 & 36.97 & 32.58 \\
2.5 & 104.28 & 87.77 & 74.56 & 63.93 & 55.33 & 48.36 \\
\bottomrule
\end{tabular}
\end{table}

\subsection{结果分析}

\parahead{温度影响最大} 题设模型中，恒温段温度由 $\SI{45}{\celsius}$ 升到 $\SI{70}{\celsius}$，烘干时长由 $\SI{67.94}{h}$ 降到 $\SI{32.58}{h}$。$\ln\tstar$ 与 $1/T_{\mathrm{set}}$ 近似成直线，斜率 $\SI{3212}{K}$，小于扩散系数中的 $\SI{3850}{K}$，因为预热平衡阶段和表面传质阻力不随扩散系数同比例变化。能量一致模型中温度的影响更强（$97.13\to\SI{35.01}{h}$，斜率 $\SI{4314}{K}$），因为温度还决定了恒速段的供热 $h(\Tair-\Twb)$。

\parahead{尺寸影响显著} 题设模型中 $\tstar\propto R_0^{1.84}$，接近内部扩散控制时的平方关系（能量一致模型为 $R_0^{1.65}$）。$\SI{50}{\celsius}$ 下半径 $\SI{1}{cm}$ 的药材只需 $\SI{16.3}{h}$，$\SI{2.5}{cm}$ 需 $\SI{87.8}{h}$。同一批药材的烘干时长由最粗者决定。

\parahead{湿度与风速的作用取决于供热是否受限} 题设模型中二者影响很小：含水浓度由 0.02 增到 0.07，烘干时长只增加 2.0\%；风速减半时长增加 13\%，加倍只减少 5\%，因为干燥主要受内部扩散控制。能量一致模型中，恒速段的干燥速率由 $h(\Tair-\Twb)$ 决定，湿度升高使湿球温度升高、风速降低使 $h$ 减小，二者影响都明显增强：含水浓度为 0.07 时烘干时长增加到 $\SI{95.68}{h}$（比 0.05 时长 39\%），0.02 时降到 $\SI{61.05}{h}$；风速减半时长增加 29\%，加倍减少 13\%。题设模型忽略蒸发吸热，会低估湿度和风速的作用。

\parahead{能耗} 以烘房相对初温的加热度时 $\int_0^{\tstar}(\Tair-\SI{28}{\celsius})\,\mathrm{d}t$ 作为烘房散热损失的代理量。恒温段温度由 $\SI{50}{\celsius}$ 升到 $\SI{60}{\celsius}$，加热度时由 $\SI{1253.5}{K.h}$ 增到 $\SI{1340.7}{K.h}$，只增加 7.0\%，而烘干时长缩短 26\%；超过 $\SI{60}{\celsius}$ 后加热度时基本不变（$\SI{70}{\celsius}$ 时为 $\SI{1347.7}{K.h}$）。

\subsection{工艺建议}

\begin{enumerate}[label=(\arabic*),leftmargin=2.2em,itemsep=1pt]
\item 恒温段温度宜取 55--$\SI{60}{\celsius}$。药典规定不宜高温烘干的药材采用低温干燥，温度一般不超过 $\SI{60}{\celsius}$\cite{chp}。在此限制内，$\SI{60}{\celsius}$ 比 $\SI{50}{\celsius}$ 缩短烘干时长 26\%（题设模型，$\SI{42.39}{h}$）至 32\%（能量一致模型，$\SI{46.56}{h}$），两种模型下都能在 2 天左右完成，而散热损失只略有增加。
\item 按粗细分级装料，或把粗根切段。半径 $\SI{2.5}{cm}$ 的药材在 $\SI{60}{\celsius}$ 下仍需 $\SI{63.9}{h}$，$\SI{1.5}{cm}$ 只需 $\SI{25.3}{h}$（表~\ref{tab:map}）。
\item 前期加强排湿和通风。恒速段的干燥速率由空气温度与湿球温度之差和换热系数决定，降低烘房湿度、提高风速能明显缩短这一时期；进入降速段后干燥受内部扩散控制，可减小风量和排湿量以节能。
\item 会收缩的药材按半径与平均含水率的关系推算。以第~\ref{sec:q4}~节的 $R=g(\Cbar)$ 计算，恒温段温度为 45、60、$\SI{70}{\celsius}$ 时，问题四药材的烘干时长分别为 $\SI{60.76}{h}$、$\SI{37.11}{h}$、$\SI{27.99}{h}$。
\end{enumerate}
